\documentclass[11pt]{article}
\usepackage[T1]{fontenc}
\usepackage{lmodern}
\usepackage[margin=1in]{geometry}
\usepackage{hyperref}
\usepackage{enumitem}
\usepackage{xcolor}
\usepackage{listings}
\usepackage{graphicx}

\definecolor{codebg}{RGB}{245,245,245}
\definecolor{codeframe}{RGB}{200,200,200}

\lstdefinestyle{commandline}{
  basicstyle=\ttfamily\small,
  backgroundcolor=\color{codebg},
  frame=single,
  rulecolor=\color{codeframe},
  breaklines=true,
  columns=flexible,
  keepspaces=true
}

\hypersetup{
  colorlinks=true,
  linkcolor=blue,
  urlcolor=blue,
  citecolor=blue
}

\title{FoldFusion User Manual}
\author{FoldFusion Development Team}
\date{\today}

\begin{document}

\maketitle

\begin{abstract}
      FoldFusion automates ligand transplantation by combining AlphaFold structures
      with experimental donor complexes. This manual explains how to install,
      configure, and operate the pipeline; how the processing stages interact; how to
      inspect results; and how to troubleshoot common issues. It provides a
      self-contained reference for researchers and engineers deploying FoldFusion on
      local workstations or compute clusters.
\end{abstract}

\tableofcontents
\newpage

\section{Quick Start}
Use these steps to execute FoldFusion exactly as shipped for the submission
deadline.

\begin{enumerate}[noitemsep]
      \item \textbf{Activate the project environment}
            \begin{lstlisting}[style=commandline]
            make venv
            source .venv/bin/activate
            make install-dev
            \end{lstlisting}
      \item \textbf{Verify configuration}: The submitted \texttt{config.toml}
            points to the analysis dataset and external binaries used for the
            final run. Confirm those paths exist on your system or adjust them
            before continuing.
      \item \textbf{Run the pipeline}: Execute the default target, which
            processes the UniProt IDs defined in \texttt{config.toml}.
            \begin{lstlisting}[style=commandline]
            make run
            # or equivalently
            foldfusion config.toml
            \end{lstlisting}
      \item \textbf{Inspect results}: Outputs are written under
            \texttt{analysis/data/foldfusion\_output}. Review
            \texttt{foldfusion\_pipeline.log} for detailed progress and
            \texttt{Evaluation/evaluation.json} for metrics.
      \item \textbf{Rebuild supporting artifacts (optional)}: Regenerate the
            analysis notebooks and figures if you need refreshed reports.
            \begin{lstlisting}[style=commandline]
            make analysis
            \end{lstlisting}
\end{enumerate}

\section{Introduction}
FoldFusion enriches AlphaFold models with experimentally observed ligands. The
pipeline stitches together external tools---DoGSite3 for pocket detection,
SIENA for structural alignment, LigandExtractor for ligand transfer, and JAMDA
for energetic refinement---and evaluates outcomes with bespoke metrics. The
project ships as a Python package with command line entry points and a managed
configuration file.

This document targets users who want to run the end-to-end workflow, adapt the
configuration, and interpret outputs. Developers extending the system will also
find an overview of its structure and testing approach.

\section{System Requirements}
\subsection{Hardware}
\begin{itemize}[noitemsep]
      \item 4 CPU cores recommended; the pipeline runs on a single core by default
            but benefits from additional concurrency.
      \item At least 16~GB RAM to accommodate SIENA database generation and JAMDA.
      \item 50~GB free storage for AlphaFold downloads, SIENA ensembles, and
            intermediate ligand files.
\end{itemize}

\subsection{Software}
\begin{itemize}[noitemsep]
      \item Linux or macOS with Python~3.12.
      \item GNU Make for the provided build targets.
      \item External binaries:
            \begin{itemize}[noitemsep]
                  \item DoGSite3 (DoGSiteScorer)
                  \item SIENA and the SIENA database generator
                  \item LigandExtractor
                  \item JAMDA scorer
            \end{itemize}
      \item Network access to the AlphaFold database if AlphaFold models are not
            cached locally.
\end{itemize}

\section{Repository Layout}
The repository is organised to keep pipeline code, tooling wrappers, and helper
scripts discoverable:
\begin{itemize}[noitemsep]
      \item \texttt{foldfusion/} --- Core package modules, including the
            orchestrator (\texttt{foldfusion.py}) and wrappers in
            \texttt{foldfusion/tools/}.
      \item \texttt{main.py} and \texttt{foldfusion/\_\_main\_\_.py} --- Entry
            points for running the pipeline via Python or console script.
      \item \texttt{config.toml} --- Central configuration file; every run reads
            from this location unless a different path is supplied.
      \item \texttt{tests/} --- Pytest-based regression and unit tests mirroring the
            package layout.
      \item \texttt{docs/}, \texttt{analysis/}, \texttt{scripts/} --- Reporting
            material, exploratory notebooks, and workflow helpers.
\end{itemize}

\section{Initial Setup}
\subsection{Environment Creation}
The project uses a virtual environment located in \texttt{.venv/}. Run:

\begin{lstlisting}[style=commandline]
make venv
source .venv/bin/activate
make install-dev
\end{lstlisting}

\noindent\textbf{Notes}
\begin{itemize}[noitemsep]
      \item \texttt{make install-dev} installs FoldFusion with developer extras,
            including Ruff, MyPy, and pytest.
      \item When deploying on shared systems, replicate the above steps inside a
            module or container that provides Python~3.12.
\end{itemize}

\subsection{External Tooling}
FoldFusion delegates heavy lifting to specialized binaries. Ensure each tool is
installed and reachable:
\begin{itemize}[noitemsep]
      \item Configure DoGSite3, SIENA, LigandExtractor, JAMDA, and (optionally) the
            SIENA database generator.
      \item Record absolute paths to these executables in \texttt{config.toml}.
      \item Verify each binary runs without interactive prompts.
\end{itemize}

\section{Configuration Reference}
The pipeline reads \texttt{config.toml} via the \texttt{Config} helper. The list
below summarises the top-level options. All paths may be absolute or relative to
the project root.

\begin{itemize}[noitemsep,leftmargin=*,labelsep=0.5em]
      \item \texttt{log\_level}: Logging verbosity (\texttt{DEBUG}, \texttt{INFO},
            \texttt{WARNING}, \texttt{ERROR}, \texttt{CRITICAL}).
      \item \texttt{log\_file}: Relative or absolute path to the log file written
            during pipeline execution.
      \item \texttt{uniprot\_ids}: Inline list of UniProt accessions to process;
            mutually exclusive with \texttt{uniprot\_ids\_file}.
      \item \texttt{uniprot\_ids\_file}: Path to a text file containing one
            UniProt ID per line.
      \item \texttt{output\_dir}: Base directory for results, logs, and
            evaluation artifacts.
      \item \texttt{pipeline\_concurrency}: Number of UniProt IDs processed in
            parallel (set to \texttt{1} for sequential execution).
      \item \texttt{resume}: When \texttt{true}, skip UniProt IDs that already
            produced a \texttt{\_SUCCESS} marker.
      \item \texttt{skip\_failed\_ids}: Skip UniProt IDs marked with
            \texttt{\_FAILED} to avoid repeating known failures.
      \item \texttt{dogsite3\_executable}: Path to the DoGSite3 binary.
      \item \texttt{dogsite3\_timeout\_seconds}: Optional timeout (seconds)
            applied to DoGSite3 runs.
      \item \texttt{dogsite3\_memory\_mb}: Optional memory ceiling in megabytes
            for DoGSite3.
      \item \texttt{siena\_db\_executable}: Path to the SIENA database builder.
      \item \texttt{siena\_db\_database\_path}: Path to an existing SIENA
            database.
      \item \texttt{pdb\_directory}: Directory holding the experimental PDB
            structures referenced by SIENA.
      \item \texttt{pdb\_format}: PDB file format flag expected by SIENA
            (\texttt{0} for \texttt{.ent.gz}, \texttt{1} for \texttt{.pdb}).
      \item \texttt{siena\_executable}: Path to the SIENA alignment executable.
      \item \texttt{siena\_max\_alignments}: Maximum number of alignments
            collected per UniProt ID.
      \item \texttt{siena\_identity}: Active-site identity threshold (0.0--1.0)
            for SIENA hits.
      \item \texttt{siena\_timeout\_seconds}, \texttt{siena\_memory\_mb}:
            Optional limits applied to SIENA executions.
      \item \texttt{siena\_db\_timeout\_seconds}, \texttt{siena\_db\_memory\_mb}:
            Controls for the database generator.
      \item \texttt{ligand\_extractor\_executable}: Path to the
            LigandExtractor binary.
      \item \texttt{ligand\_extractor\_timeout\_seconds},
            \texttt{ligand\_extractor\_memory\_mb}: Execution limits for ligand
            extraction.
      \item \texttt{jamda\_scorer\_executable}: Path to the JAMDA CLI.
      \item \texttt{jamda\_timeout\_seconds}, \texttt{jamda\_memory\_mb}:
            Optional limits for JAMDA optimisation.
\end{itemize}

\subsection{Managing UniProt Inputs}
Choose exactly one of \texttt{uniprot\_ids} or \texttt{uniprot\_ids\_file}. The
configuration loader raises a \texttt{ValueError} if both are set. For batch
processing, prefer the file-based option to keep \texttt{config.toml} stable.

\begin{lstlisting}[style=commandline]
# Example snippet: inline UniProt list
uniprot_ids = ["Q8CA95", "Q9QYJ6", "P12345"]
# output_dir adjusts where Results/, Logs/, and Evaluation/ are created
output_dir = "/data/foldfusion/runs/run_2025_05_01"
\end{lstlisting}

\subsection{External Tool Paths}
Record absolute paths when possible to avoid ambiguity. When sharing
configuration files across machines, document machine-specific changes in
project notes instead of committing them.

\section{Running the Pipeline}
\subsection{Make Targets}
Convenience targets coordinate common tasks:
\begin{itemize}[noitemsep]
      \item \texttt{make run} --- Executes the pipeline using \texttt{config.toml}.
      \item \texttt{make test} --- Runs the pytest suite with coverage reporting.
      \item \texttt{make lint} --- Executes Ruff linting and MyPy type checking.
      \item \texttt{make format} / \texttt{make format-fix} --- Apply Ruff
            formatting (with optional autofix).
      \item \texttt{make analysis} --- Rebuild analysis notebooks, tables, and
            figures under \texttt{analysis/}.
      \item \texttt{make docs} or \texttt{make thesis} --- Build LaTeX outputs in
            \texttt{docs/} or \texttt{thesis/} respectively.
      \item \texttt{make docs-clean} --- Remove previously generated LaTeX
            artefacts (e.g. \texttt{*.aux}, \texttt{*.log}, \texttt{*.toc}) before
            rebuilding documentation targets.
\end{itemize}

\subsection{Command Line Interfaces}
The package exposes two equivalent CLI entry points:
\begin{itemize}[noitemsep]
      \item \texttt{python main.py} --- Runs FoldFusion with \texttt{config.toml} in
            the project root.
      \item \texttt{foldfusion [CONFIG]} --- Console script installed by
            \texttt{pip}. The configuration path is optional and defaults to
            \texttt{config.toml} in the current directory.
\end{itemize}

\begin{lstlisting}[style=commandline]
# Run with explicit configuration file
foldfusion /path/to/custom_config.toml
\end{lstlisting}

FoldFusion writes progress logs to the configured file and mirrors key events to
stdout. Long-running steps (notably SIENA database generation) emit periodic
messages so you can monitor progress.

\section{Pipeline Workflow}
Each UniProt accession flows through six stages handled by dedicated helper
classes inside \texttt{foldfusion/tools/}:
\begin{enumerate}[noitemsep]
      \item \textbf{AlphaFold retrieval} --- \texttt{AlphaFoldFetcher} downloads the
            latest model (preferring v4) and trims segments with low pLDDT scores.
      \item \textbf{Pocket detection} --- \texttt{Dogsite3} runs DoGSiteScorer to
            identify the highest-scoring binding site. The resulting EDF file becomes
            the template for structural alignment.
      \item \textbf{Structural alignment} --- \texttt{Siena} aligns the pocket
            against experimental complexes in the SIENA database and ranks hits by
            active-site identity and RMSD.
      \item \textbf{Ligand transplantation} --- \texttt{LigandExtractor} copies
            ligands from the SIENA ensembles into the AlphaFold frame, emitting SDF
            files organised by donor PDB code.
      \item \textbf{Refinement} --- \texttt{JamdaScorer} optimizes each transplanted
            ligand using energy minimisation and records both the refined structure
            and the donor reference.
      \item \textbf{Evaluation} --- \texttt{Evaluator} calculates local RMSD,
            Transplant Clash Score (TCS), and Ligand Environment Verification (LEV)
            metrics for pre- and post-JAMDA stages.
\end{enumerate}

The orchestrator coordinates these components, tracks duration, handles error
recovery, and writes stage-specific logs for traceability.

\section{Output Structure}
FoldFusion mirrors outputs under the configured \texttt{output\_dir}.
\begin{itemize}[noitemsep]
      \item \texttt{Results/\textless UniProt\textgreater/AlphaFold/} --- Downloaded
            and processed AlphaFold PDB files.
      \item \texttt{Results/\textless UniProt\textgreater/Dogsite3/} --- EDF pocket
            descriptors from DoGSite3.
      \item \texttt{Results/\textless UniProt\textgreater/Siena/} --- SIENA log
            files, \texttt{resultStatistic.csv}, and \texttt{ensemble/} PDBs for
            each alignment.
      \item \texttt{Results/\textless UniProt\textgreater/LigandExtractor/} --- SDF
            ligands grouped by donor PDB code plus \texttt{ligand\_structure.json}.
      \item \texttt{Results/\textless UniProt\textgreater/JamdaScorer/} --- Refined
            ligand SDFs that preserve a \texttt{ref\_path} back to the donor file.
      \item \texttt{Results/\textless UniProt\textgreater/\_SUCCESS} or
            \texttt{\_FAILED} --- Markers controlling resume logic.
      \item \texttt{Evaluation/evaluation.json} --- Aggregated metrics keyed by
            UniProt ID, donor PDB code, ligand identifier, and pipeline stage.
      \item \texttt{foldfusion\_pipeline.log} --- Log file reflecting the level
            specified in \texttt{config.toml}.
\end{itemize}

\section{Monitoring and Logging}
\begin{itemize}[noitemsep]
      \item Logs capture timestamps, stage transitions, command invocations, and
            error details. Adjust \texttt{log\_level} to control verbosity.
      \item Each tool wrapper records the exact command line before execution; use
            these entries to rerun failing commands manually.
      \item Errors triggered by optional steps (e.g., ligand extraction) are logged
            and skipped so downstream tasks can proceed where possible.
\end{itemize}

\section{Testing and Quality Assurance}
\begin{itemize}[noitemsep]
      \item \texttt{pytest} suites live under \texttt{tests/}. Run
            \lstinline[style=commandline]|pytest --cov=foldfusion --cov-report=term-missing|
            to inspect coverage.
      \item Mock external binaries in tests to keep them deterministic and fast.
      \item Type hints cover all new code; verify with \texttt{mypy foldfusion/}.
      \item Before submitting changes, execute \texttt{make lint} and
            \texttt{make test} and record the outcomes.
\end{itemize}

\section{Operational Tips}
\subsection{Managing Concurrency}
Increase \texttt{pipeline\_concurrency} once external tools and the filesystem
can handle parallel workloads. Monitor memory usage when running multiple SIENA
jobs simultaneously.

\subsection{Caching Resources}
\begin{itemize}[noitemsep]
      \item AlphaFold models download per UniProt ID; reuse the generated
            \texttt{AlphaFold/} directories to avoid repeated transfers.
      \item The SIENA database generator writes the database into the directory
            defined by \texttt{siena\_db\_database\_path}. Rebuild only when the
            source PDB archive changes.
\end{itemize}

\subsection{Extending the Pipeline}
\begin{itemize}[noitemsep]
      \item New stages should follow the pattern established by existing tool
            wrappers: inherit from \texttt{Tool}, expose a \texttt{run()} method, and
            persist outputs in a dedicated subdirectory.
      \item Update \texttt{foldfusion/foldfusion.py} to insert additional stages and
            extend the evaluation hooks if new metrics are required.
\end{itemize}

\section{Troubleshooting}
\subsection{Common Issues}
\begin{itemize}[noitemsep]
      \item \textbf{AlphaFold download failures}: Verify network connectivity and
            confirm the UniProt accession exists in the AlphaFold database. The
            fetcher falls back from model v4 to v3 automatically.
      \item \textbf{Missing EDF files}: DoGSite3 must complete successfully before
            SIENA runs. Check \texttt{Dogsite3/} for error messages.
      \item \textbf{SIENA database errors}: Ensure the database path exists or that
            the generator binary has write permission to the target directory.
      \item \textbf{Ligand extraction skips}: Warnings about missing SDF files often
            indicate that the donor PDB lacks ligands for the selected chain. Review
            \texttt{ligand\_structure.json} for details.
      \item \textbf{JAMDA crashes}: Large ligands may exceed memory limits. Increase
            \texttt{jamda\_memory\_mb} or disable the limit temporarily.
      \item \textbf{Evaluation gaps}: If evaluations show \texttt{null} metrics,
            confirm that both the ensemble path and ligand path were produced. The
            evaluator skips metrics when prerequisites are missing.
\end{itemize}

\subsection{Resume and Recovery}
When \texttt{resume = true}, FoldFusion looks for \texttt{\_SUCCESS} markers in
per-UniProt result directories. Delete the marker to force a rerun for that
protein. Failures leave a \texttt{\_FAILED} file containing the captured error
message; adjust configuration and remove the file to retry the affected ID.

\section{Glossary}
\begin{itemize}[noitemsep]
      \item \textbf{AlphaFold model}: Predicted protein structure used as the host
            for ligand transplantation.
      \item \textbf{EDF (Pocket Description File)}: Output from DoGSite3 describing
            binding pocket geometry.
      \item \textbf{Ensemble}: SIENA-generated PDB file aligning the donor complex
            to the AlphaFold structure.
      \item \textbf{Local RMSD}: RMSD computed over residues within 6~\AA\ of the
            ligand position after alignment.
      \item \textbf{TCS (Transplant Clash Score)}: FoldFusion metric assessing steric
            clashes between transplanted ligand and the receptor.
      \item \textbf{LEV (Ligand Environment Verification)}: Optional comparison of
            ligand environments when donor and target share perfect active-site
            identity.
\end{itemize}

\end{document}
